\begin{abstract}
We refute Kaplansky's quasitrace conjecture by constructing a separable
unital complex \(C^*\)-algebra that admits normalized \(2\)-quasitraces
but has no tracial state. As a consequence, we show that the minimal tensor
product of two unital simple stably finite complex \(C^*\)-algebras can be
properly infinite, even when one factor is the reduced free-group algebra
\(C_r^*(\mathbb F_2)\).
\end{abstract}

\section{Introduction}
\label{sec:introduction}

A unital \(C^*\)-algebra is \emph{stably finite} when every matrix
algebra over it has no nonunitary isometry, and is \emph{properly
infinite} when it contains two isometries with orthogonal ranges.
Milh{\o}j and R{\o}rdam asked whether the minimal tensor product of two
unital simple stably finite \(C^*\)-algebras must again be stably finite
\cite[Question~2.5]{MR}. They proved that this question is equivalent
to the quasitrace problem and that a negative answer to that problem would yield a
properly infinite tensor product with \(C_r^*(\mathbb F_\infty)\)
\cite[Section~2 and Theorem~3.1]{MR}. We construct the required
quasitrace counterexample. The resulting tensor product can be properly
infinite with one factor equal to the reduced \(C^*\)-algebra
\(C_r^*(\mathbb F_2)\) of the free group on two generators.

A quasitrace retains positivity and the identity
\(\tau(x^*x)=\tau(xx^*)\), but is required to be linear only on
commutative subalgebras, together with compatibility with real and
imaginary parts. A \(2\)-quasitrace must also extend to the two-by-two
matrix algebra. Whether this extra condition forces full linearity
is the general quasitrace problem. It is equivalent to Kaplansky's
question whether every \(AW^*\)-factor of type \(\mathrm{II}_1\) is a
von Neumann algebra; see \cite[Section~2]{MR} and
\cite[Conjectures~1.1--1.2 and Theorem~4.11]{Gow}.
Here the term \emph{von Neumann algebra} has its abstract \(W^*\)
meaning: the \(C^*\)-algebra admits a Banach-space predual.

Gow's historical account explains that Kaplansky introduced
\(AW^*\)-algebras in 1951 to characterize \(W^*\)-algebras by intrinsic
algebraic conditions, independent of a chosen Hilbert-space
representation. The factor types have different outcomes: every type
\(\mathrm{I}\) \(AW^*\)-factor is a \(W^*\)-algebra, whereas type
\(\mathrm{III}\) counterexamples exist. The type \(\mathrm{II}_1\)
factor question is the finite case addressed here
\cite[Introduction]{Gow}.

Quasitraces also connect the projection-theoretic origins of the question
to dimension theory for general \(C^*\)-algebras. Cuntz developed
dimension functions for simple algebras \cite{Cuntz1978}, and
Blackadar and Handelman established the correspondence between
\(2\)-quasitraces and lower semicontinuous dimension functions
\cite{BH}. Haagerup credits Blackadar and Handelman with proving that
a positive answer to Kaplansky's factor question would imply that every
\(2\)-quasitrace on a unital \(C^*\)-algebra is a trace. He also recalls
their realization of such quasitraces as pullbacks of quasitraces on
finite \(AW^*\)-algebras along unital \(*\)-homomorphisms
\cite[Introduction and Proposition~3.3]{Haagerup}.
Blackadar and R{\o}rdam later developed the extension of states on
preordered semigroups that underlies the general existence criterion
\cite{BR}.

The exact case was settled by Haagerup: every \(2\)-quasitrace on a
unital exact \(C^*\)-algebra is a trace \cite[Theorem~5.11]{Haagerup}.
A first draft was completed in 1991 and circulated thereafter; the
paper was published in 2014. Without
exactness, the existence of quasitraces is governed by a different
condition. The Cuntz--Blackadar--Handelman existence theorem provides a
normalized \(2\)-quasitrace whenever no matrix algebra over the given
unital algebra is properly infinite \cite[Theorem~2.3]{MR}.
Thus a unital algebra with that property but without tracial states
would answer the general question negatively.

We construct such an algebra and obtain a fixed positive pair
witnessing nonadditivity. Throughout, a normalized \(2\)-quasitrace
has value one on the unit of its domain; its corner-preserving
extension to \(M_2\) has value two on the matrix identity.
The definitions are recalled in Section~\ref{sec:preliminaries}.

\Needspace{10\baselineskip}
\begin{theorem}
\label{thm:main}
There exist a separable unital complex \(C^*\)-algebra \(A\), admitting
a normalized \(2\)-quasitrace, and positive contractions \(a,b\in A\)
such that every normalized \(2\)-quasitrace \(\tau\) on \(A\) satisfies
\[
 \tau(a+b)-\tau(a)-\tau(b)\ge\frac1{144}.
\]
\end{theorem}

Theorem~\ref{thm:main} refutes the conjecture that every normalized
\(2\)-quasitrace on every unital complex \(C^*\)-algebra is a tracial
state. By the equivalence recalled above, it also gives a negative
answer to Kaplansky's \(\mathrm{II}_1\) \(AW^*\)-factor question.
The constructed algebra is necessarily nonexact, by Haagerup's
result.

The matrix-extension requirement is essential to this conclusion.
Haagerup reports Kirchberg's 2006 unpublished example of a
\(1\)-quasitrace that is not a \(2\)-quasitrace, and hence is not a trace
\cite[Introduction]{Haagerup}. This example lacks the matrix extension
required here.
Gow's recent work places the established factor--quasitrace equivalence
alongside further equivalent formulations involving maximal abelian
subalgebras and bidual retractions \cite[Theorem~4.11]{Gow}.

The tensor-product conclusion and two related consequences can be stated
together. All tensor products below are minimal \(C^*\)-tensor products.

\Needspace{12\baselineskip}
\begin{corollary}[Tensor products and stable finiteness]
\label{cor:stable-finiteness}
\leavevmode
\begin{enumerate}[label=\textup{(\roman*)}]
\item There are two unital simple stably finite complex \(C^*\)-algebras
whose minimal tensor product is properly infinite. One factor can be
taken to be \(C_r^*(\mathbb F_2)\).
\item There is a separable unital complex \(C^*\)-algebra \(B\) admitting
a normalized \(2\)-quasitrace such that
\(B\otimes_{\min}C_r^*(\mathbb F_2)\) is properly infinite and admits no
normalized \(2\)-quasitrace.
\item There is a separable unital stably finite complex \(C^*\)-algebra
with no tracial state.
\end{enumerate}
\end{corollary}

The three assertions use related but different algebras. The simple
tensor-product example comes from a non-von-Neumann \(AW^*\)-factor of
type \(\mathrm{II}_1\), with no separability assertion for that factor.
The quasitrace-loss example uses the explicit separable algebra \(B\)
constructed below; we prove that no matrix algebra over \(B\) is properly
infinite, not that \(B\) is stably finite. A finitely generated subalgebra
of the \(AW^*\)-factor gives the separable stably finite traceless example,
without a simplicity assertion. The short deductions from established
trace-obstruction and tensor-product results are given in
Section~\ref{sec:consequences}.

\subsection*{The construction}

We first build a unital algebra \(B=C^*(1,x_1,\ldots,x_{24})\) such that
\begin{equation}
 \sum_{i=1}^{24}x_i^*x_i=1,\qquad
 \sum_{i=1}^{24}x_ix_i^*\le\tfrac23\,1,
 \label{eq:intro-relations}
\end{equation}
while no \(M_m(B)\), \(m\ge1\), contains two isometries with orthogonal
ranges. A tracial state would make the two sums in
\eqref{eq:intro-relations} have the same value, so \(B\) has none.
These relations instantiate Haagerup's characterization of the absence
of tracial states by a finite tuple with column sum one and a strict
row bound \cite[Lemma~2.1]{Haagerup}.
Milh{\o}j and R{\o}rdam proposed separating a uniformly bounded length
of such a tuple from increasingly delayed matrix proper infiniteness
\cite[Remark~4.7]{MR}. The construction below uses a bounded product
and a generated subalgebra, with a separate test for every matrix size
and polynomial degree. The fixed number of generators preserves the
trace obstruction while all finite polynomial candidates are excluded.

The generators are shifts between finite-rank vector bundles. For
each pair of positive integers \((m,R)\), we choose a compact base
\(Y_{m,R}\) and bundles \(E_j\to Y_{m,R}\) for all \(j\geq0\), with
\(E_0\) trivial. Only levels \(0,\ldots,R\) enter the finite topological
test. A polynomial matrix \(P\in M_{m,2m}(B)\) whose entries contain
only words of length at most \(R\) restricts to a bundle map
\[
 P_0:E_0^{\oplus2m}\longrightarrow
 \left(\bigoplus_{j=0}^{R}E_j\right)^{\oplus m}.
\]
The target's Chern class forbids the trivial image bundle that would
result from \(P^*P\geq\tfrac12 I\). Write \(S_i^{m,R}(y)\) for the
corresponding forward shift on the Hilbert direct sum
\(\bigoplus_{j\geq0}E_j(y)\). Uniform norm bounds allow the simultaneous definitions
\[
 x_i=\bigl(S_i^{m,R}(y)\bigr)_{m,R\geq1,\;y\in Y_{m,R}},
 \qquad B=C^*(1,x_1,\ldots,x_{24}),
\]
in one bounded operator product. The number of generators and the row
bound in \eqref{eq:intro-relations} are uniform, while the bundle ranks
and base dimensions may depend on \((m,R)\). A polynomial matrix
approximating a candidate pair of isometries has the same coefficients
at every coordinate. Its finite degree selects one of the tests already
present in the product, which obstructs that approximation.

The topological obstruction follows the use of characteristic classes
in the constructions of Villadsen \cite[Section~3]{Villadsen1999}
and R{\o}rdam \cite[Proposition~4.5]{Rordam2003}.
Here projective-space factors retain a nonzero Chern class after a
controlled trivial summand is allowed. The parameters are selected
anew for each finite test.

Two ingredients keep the expansion and the topological obstruction
compatible. First, a family of bundle maps whose collective image
expands every source subspace by a factor of two can be normalized
continuously to have column sum one and row sum at most \(2/3\).
The determinant-potential approach is related to operator scaling
\cite[Section~4]{Gurvits2003} and to geodesically convex matrix
optimization \cite[Corollary~2.11]{SraHosseini2015}.
We prove strict convexity and existence of a unique minimizer for the
present rectangular maps; uniqueness makes the normalization compatible
with changes of bundle frame.
Second, the bundle tests separate small and large source subspaces:
sixteen maps into reduced copies of old line bundles handle the
small ones, while eight generic maps into a new line-bundle block
handle the large ones. This keeps the generator count fixed as the
matrix size and word-length bound increase.

Finally, set \(A=M_{48}(B)\) and apply Theorem~\ref{thm:existence}
directly to \(A\). A finite telescoping construction, using only
linearity on commuting elements, gives the pair in
Theorem~\ref{thm:main}. This last step works for every normalized
\(1\)-quasitrace on \(A\), hence for every normalized \(2\)-quasitrace
there, and requires no formula for a chosen quasitrace.

\subsection*{Organization}

Section~\ref{sec:preliminaries} records the operator-algebra
conventions and the existence theorem. Section~\ref{sec:bundle-map-lemmas}
proves the incidence and scaling lemmas. The finite bundle tests
are assembled in Section~\ref{sec:bundle-tests}. Section~\ref{sec:algebra}
constructs \(B\), excludes proper infiniteness in every matrix size,
exhibits the positive pair, and completes the proof of
Theorem~\ref{thm:main}. Section~\ref{sec:consequences} proves
Corollary~\ref{cor:stable-finiteness}.
